\begin{textsample}{POS Dim 3 – es – Score 15.00 – es/es\_inf\_12\_camp\_exp\_escuta\_fala.txt}  \label{ex:f3_pos_067}
Campo de Experiências : \textbf{Escuta}, Fala, Pensamento e Imaginação A constituição do pensamento está associada ao domínio e à apropriação da linguagem verbal.
Os \textbf{bebês} e as \textbf{crianças}, à medida que vão se apropriando dos sentidos e dominando a língua materna, vão ampliando e enriquecendo \textbf{progressivamente} seu vocabulário e utilizando recursos de expressão e de compreensão cada vez mais complexos, tornando a língua seu veículo privilegiado de interação.
É no convívio e nas interações com seus pares que elas aprendem a falar, a ouvir, a compreender seu contexto, vivenciando experiências que potencializam sua participação na cultura.
Sorrir, falar, \textbf{imitar}, tagarelar, \textbf{inventar} histórias, fazer perguntas, expressar suas ideias e opiniões, defender seus pontos de vista são capacidades que vão sendo desenvolvidas pelas \textbf{crianças} e que marcam significativas experiências no campo Escuta, Fala, Pensamento e Imaginação.
Toda esta experiência com a linguagem verbal dada pela \textbf{criança}, além de ser ampliada gradativamente, conforme vai se desenvolvendo, dialogam com outras linguagens, como o pensamento ( sobre si, o outro, o mundo e a língua ) e a imaginação das \textbf{crianças}.
Escutar e falar são atos que estão intrinsicamente ligados e constituem a língua e o pensamento humanos desde o nascimento.
Ao ingressarem na Educação \textbf{Infantil}, os momentos de escutar e falar das \textbf{crianças} são atos transversais que perpassam todos os campos de experiências.
As DCNEI ( Parecer CNE/CEB nº20/09 ) trazem que as práticas pedagógicas na Educação \textbf{Infantil} devem garantir experiências que “ favoreçam a imersão das \textbf{crianças} nas diferentes linguagens e o progressivo domínio por elas de diferentes gêneros e formas de expressão : gestual, verbal, \textbf{plástica}, \textbf{dramática} e musical ”.
Assim, este campo de experiências que evidencia muito claramente a linguagem verbal não se separa completamente das outras linguagens expressas nos outros campos : corporal, musical, \textbf{plástica} e \textbf{dramática}.
A gestualidade, movimento exigido nas \textbf{brincadeiras} ou jogos corporais, a aquisição da linguagem verbal ( oral e escrita ), ou em libras, potencializam tanto a comunicação quanto a organização do pensamento das \textbf{crianças} e sua participação na cultura.
Com efeito, ao oportunizar aos \textbf{bebês} e às \textbf{crianças}, nos espaços da \textbf{creche} e da \textbf{pré-escola}, a \textbf{escuta} de histórias, a participação nas conversas, nas descrições, nas narrativas elaboradas individualmente ou em grupo, nas implicações com as múltiplas linguagens, a \textbf{criança} se constitui ativamente como sujeito singular, pertencente a um grupo social.
Além disso, o contato com histórias, contos, canções, rimas, leitura de imagens ; o contato com as letras ; identificação de palavras, fábulas, poemas, cordéis, livros de diferentes gêneros literários ; a \textbf{escuta} e dramatização de histórias ; a participação na produção de textos escritos ; dentre outros ; propiciam a construção de novos conhecimentos a respeito da linguagem verbal, desenvolvem o gosto pela leitura, estimulam à imaginação, ampliam o conhecimento de mundo, além de promoverem a apropriação de novos \textbf{gestos}, falas, histórias e escritas ( convencionais ou não ).
O papel do professor/a é o de articular as experiências e saberes \textbf{infantis}, pois ele/a é o propositor de atividades, que envolvem múltiplas linguagens possibilitando às \textbf{crianças} explorar e dar sentido ao mundo.
Na interação com seus pares e os \textbf{adultos} no ambiente escolar, as \textbf{crianças} buscam captar os signos e símbolos socialmente construídos, presentes nos comportamentos dos parceiros humanos, expressos na oralidade e na escrita presentes no ambiente escolar, levando -as a, \textbf{progressivamente}, entender as formas de comunicação, investigar e reconhecer os diferentes usos sociais da linguagem verbal, imergindo na cultura do escrito.
Com relação à linguagem escrita, é no convívio com textos escritos que as \textbf{crianças} vão construindo hipóteses sobre a escrita que se revelam, inicialmente, em rabiscos e garatujas e, à medida que vão progredindo em seu desenvolvimento e conhecendo as letras, em escritas espontâneas, não convencionais, mas já indicativas da compreensão da escrita como sistema de representação da língua, que vão sendo emergidas no universo da linguagem oral e escrita.
As DCNEI’s reconhecem esta linguagem como de interesse pela \textbf{criança} desde cedo.
Também chamam a atenção do professor para que suas práticas pedagógicas para tal linguagem sejam coerentes com o que se conhece como sendo especificidades da primeira infância.

% matched lemmas: adulto, bebê, brincadeira, creche, criança, dramático, escuta, gesto, imitar, infantil, inventar, plástico, progressivamente, pré-escola
\end{textsample}
